양자 위협에 대응하여 제안된 격자 기반 Fiat--Shamir 서명 스킴 \texttt{HAETAE}는 수학적 안전성과
별개로 실제 구현 환경에서 오류주입공격(Fault Injection Attack, FIA)에 취약할 수 있다. 특히 서명의
응답 연산 $\zvec=\yvec+(-1)^b c\svec_1$은 비밀키와 일회용 난수가 결합되는 유일한 지점으로,
이 연산의 무결성이 보장되지 않으면 비밀 정보가 직접 노출될 수 있다. 본 논문에서는 \texttt{HAETAE}의
서명 구조를 분석하여 응답 연산에 대한 세 가지 오류 지점($c\svec$ 스킵, $+\yvec$ 스킵, 거부 우회)을
도출하고, 그 중 $+\yvec$ 덧셈을 건너뛰는 경우 단 한 번의 오류만으로 비밀키 $\svec_1$을 직접
복원할 수 있음을 보인다. 전체 서명 구현에서
(ARM 기반 MCU에 에뮬레이션된 단일 오류 모델 하에) 오류가 주입된 단일 서명으로부터 $\svec_1$의 768개 계수 전부를
100\% 복원함을 확인하였으며, 복구 파이프라인의 정확성은 자체 검증으로 입증하였다. 또한 본
논문은 연산 재계산 검증, 서명 경계 노름 재검사, 비밀키 무결성 검사를
비교(분기) 연산에 의존하지 않는 감염형 마스킹으로 통합한 경량 대응 기법
(IRV, Infective Randomized Verification, 감염형 무작위 검증)를
제안한다. 제안 기법은 (i)~거부 우회를 포함한 7개 오류 지점을 모두 차단하고, (ii)~탐지 분기를
건너뛰는 2차 오류에 내성을 가지며, (iii)~전체 이중 연산 대비 절반 수준의 서명 시간으로
동작함을 보인다.
